% 摘要与关键词（共享内容，供 main-zh.tex 和 main-jos.tex 引用）
% 修改摘要只需编辑此文件，两个入口文件自动同步。

% ========== 中文摘要 ==========
\newcommand{\AbstractContentZh}{%
微服务架构下，日志来源高度分散、格式各异，总量随并发请求线性增长，节点资源、网络带宽和集中存储的压力随之显著增加。传统全量采集模式难以兼顾高价值日志保留与系统开销控制。本文提出一种基于动态关注清单的微服务日志定向采集方法，核心思路是``只采必要日志''。主要贡献包括四个方面：（1）构建网关预筛选、节点定向采集与中心二次过滤的三层协同架构，以 Sidecar 模式部署，从源头实现采集前端减量；（2）提出动态关注度评分模型（DASM），以频次、错误率、延迟严重度和热度趋势四个归一化因子加权求和，结合指数时间衰减和负载自适应权重调节，动态生成关注清单，无需静态配置或业务代码侵入即可实时识别高价值 URL 模式，并证明评分有界性和退化兼容性；（3）提出定向策略三次转换算法，通过统一的 URL 泛化函数和清单版本控制维持跨层语义一致；（4）引入固定缓存块与压力感知指数退避机制，辅以 BoltDB 本地兜底，保障资源受限环境下的可靠传输。原型基于 Go 和 Grafana Loki 实现。八节点集群对比实验（180\,s、并发 50、$n{=}10$ 配对）显示，定向采集将 Loki 入库量从 4392 条降至 72 条（降幅 98.4\%，$p<0.01$，Cohen's $d \approx 378$），Agent CPU 与内存开销分别降低 28.2\% 和 2.3\%。DSB-Lite 真实部署（7 个核心微服务）的 $n{=}5$ 配对实验进一步验证了方法在多服务真实场景下的有效性（减量 $25.9\% \pm 0.9\%$，$p<0.001$）。十六/三十二节点规模复核和 1--1000 节点扩展性外推进一步验证了方法的可扩展性。参数敏感性分析与权重寻优（三场景 60 组合网格搜索）表明，Top-$K$ 和延迟阈值 $T$ 对入库量和覆盖率有显著影响，自适应权重机制相对默认配置可提升 8--15\% 减量率，为生产部署提供了调优依据。该方法在显著降低日志入库量的同时有效保留了异常与慢请求的关键上下文，为微服务可观测性提供了高效、低侵入的采集前端解决方案。%
}

% ========== 英文摘要 ==========
\newcommand{\AbstractContentEn}{%
In microservice architectures, log sources are highly distributed with heterogeneous formats, and their aggregate volume grows linearly with request concurrency, imposing significant overhead on node resources, network bandwidth, and centralized storage. Traditional full-volume collection struggles to balance high-value log retention against system overhead. This paper proposes a dynamic attention list-based directed log collection method for microservices, guided by the principle of ``collect only what matters.'' The main contributions span four aspects: (1) a three-layer collaborative architecture comprising gateway pre-filtering, node-side directed collection, and center-side secondary filtering, deployed in Sidecar mode to achieve volume reduction at the collection frontend; (2) a Dynamic Attention Scoring Model (DASM) that combines four normalized factors---frequency, error rate, latency severity, and trend---with exponential time decay and workload-adaptive weight adjustment to dynamically generate attention lists, with proven score boundedness and degradation compatibility; (3) a triple strategy transformation algorithm that maintains cross-layer policy semantic consistency through a unified URL generalization function and list version control; (4) fixed-size cache blocks coupled with pressure-aware exponential backoff and BoltDB local fallback to ensure reliable transmission under resource-constrained environments. A prototype built with Go and Grafana Loki is evaluated on a WSL/Docker cluster. Eight-node comparison experiments (180\,s, concurrency 50, $n{=}10$ paired) show that directed collection reduces Loki-ingested logs from 4392 to 72 entries (98.4\% reduction; $p<0.01$, Cohen's $d \approx 378$), while agent CPU and memory overhead decrease by 28.2\% and 2.3\%, respectively. A $n{=}5$ paired experiment on a real DSB-Lite deployment (7 core microservices) further validates the method's effectiveness in multi-service real scenarios ($25.9\% \pm 0.9\%$ reduction, $p<0.001$). Scale validations at 16 and 32 nodes and 1--1000 node scalability extrapolation further confirm the method's scalability. Parameter sensitivity analysis and weight optimization (60-combination grid search across three scenarios) reveal that Top-$K$ and latency threshold $T$ significantly impact ingestion volume and coverage, with the adaptive weight mechanism providing an additional 8--15\% reduction over the default configuration, providing tuning guidance for production deployment. The method substantially reduces log ingestion volume while effectively retaining critical context for anomalous and slow requests, offering an efficient, low-intrusion collection-frontend solution for microservice observability.%
}

% ========== 中文关键词 ==========
\newcommand{\KeywordsZh}{微服务;分布式日志采集;动态关注度评分;定向过滤;参数敏感性;Grafana Loki}

% ========== 英文关键词 ==========
\newcommand{\KeywordsEn}{microservice; distributed directed log collection; dynamic attention scoring; directed filtering; parameter sensitivity; Grafana Loki}
